

\chapter{Methodology}\label{chap:methodology}

The implementation treats the open-source LLaMEA library as the orchestration engine, focusing our engineering efforts on a highly defensive evaluation domain.

\begin{itemize}
    \item \textbf{Stage 1: Dynamic Prompting \& Alteration:} The system initializes a clean \texttt{ioh} BBOB landscape and immediately wraps it in an additive Gaussian noise layer. The LLM is then prompted dynamically with the specific BBOB Problem ID, forcing it to specialize its strategy.
    
    \item \textbf{Stage 2: The \texttt{evaluate()} Sandbox:} The core engineering contribution. The pipeline extracts the LLM's raw Python generation and compiles it (using \texttt{exec}) within a strictly controlled sandbox.
    
    \item \textbf{Stage 3: Robustness \& Crash Defenses:} To prevent ``lucky'' evaluations caused by landscape noise, the algorithm is forced to run 5 independent times. The mean absolute error across these runs determines the fitness score. Infinite loops are caught via a 10-second \texttt{func\_timeout}, and runtime exceptions generate Python Tracebacks that are fed directly back to the LLM to facilitate self-debugging.
    
    \item \textbf{Stage 4: Evolution \& Artifact Generation:} LLaMEA natively handles a 1+1 Evolution loop for 30 generations per problem. Upon completion, the pipeline logs the generation metadata and saves the final champion solver as an independent Python artifact.
\end{itemize}